% ============================================================
% REWRITE DRAFT — Chapter 3 in full, economics-paper register. Five sections.
% Every estimate is from thesis/tables/ch3_estimates.csv (Stata, HC3, the
% documented control set incl. years_since_ofsted, primary spec n=99). The
% five-outcome tables, entry rates, gaps, typology detail, pseudo-P8 check
% and stability evidence stay in Appendix 3A. Numbers convert to macros at
% integration.
% ============================================================

\section{Introduction}
\label{sec:p2_intro}

Does the culture of a school affect how much its pupils learn? This paper
estimates the association between the warmth and strictness of English
secondary schools, measured by direct observation, and the progress their
pupils make between the end of primary school and GCSE. It uses the enacted
culture scores constructed in the previous chapter for \NTierOne{} visited
schools, and the Progress~8 measure published by the Department for
Education, which records each school's average pupil progress relative to
pupils nationally with the same prior attainment.

The design is an association with several features that strengthen it short
of identification. The outcome is value-added, so it already conditions on
where pupils started; the culture measures come from documents and visits that
mostly predate or coincide with the outcome cohort; the regressions condition
on intake composition, prior attainment, school type, size, location and the
school's pre-pandemic inspection grade; and the sample excludes schools that
admit pupils only at 13 or 14, whose Progress~8 attributes to them progress
they cannot have caused. What the design cannot do is rule out that some
persistent characteristic of a school produces both its culture and its
results, and the estimates are presented as associations throughout.

Three results follow. First, both dimensions predict progress, with about
equal weight: a one standard deviation increase in observed warmth is
associated with $0.14$ of a grade higher Progress~8, and in observed
strictness with $0.13$, both significant at the one per cent level, and both
stable across every treatment of the inspection-grade control. Entered alone
each dimension carries a coefficient of about $0.22$, so the two share some
variance but each carries independent signal. Second, the association does not
run through observed teaching quality: adding the visit's teaching score
leaves both culture coefficients unchanged and the teaching coefficient
itself is close to zero. Third, the schools that combine high warmth with high
strictness outperform the rest by a third of a grade, while the interaction
between the two dimensions is null; the dimensions add rather than amplify
one another. Substituting the headteacher's own account of the school for the
observed scores removes the warmth association entirely, and the strictness
score read from inspection reports reproduces the strictness result across
the national population with the inspection grade held constant.

The paper proceeds as follows. Section~2 describes the outcome, the culture
measures, the controls and the sample. Section~3 sets out the specification.
Section~4 reports the results and their robustness. Section~5 concludes.

% ============================================================
\section{Data}
\label{sec:p2_data}

\textbf{Outcomes.} The outcome is Progress~8, the Department for Education's
school-level value-added measure: each pupil's attainment across eight GCSE
slots is compared with the national average for pupils with the same Key
Stage~2 starting point, and the differences are averaged across the school. A
school at zero is exactly average; $+0.5$ means half a grade more per subject
than its intake would predict. The paper reports overall Progress~8 and its
English and Maths components, which are the cleanest measures the tables offer
because every pupil sits both examinations. The remaining components, EBaC and
Open, are partly determined by which subjects a school enters pupils for, and
are reported in Appendix~3A together with a direct analysis of subject entry
rates and a note on how the components are constructed. The primary outcome is
the two-year average of the 2022/23 and 2023/24 cohorts, the two post-pandemic
years for which a Key Stage~2 baseline exists; the 2024/25 cohort has no
baseline, and Appendix~3A reports a validated pseudo-progress measure for it.

\textbf{Culture measures.} The explanatory variables are the enacted warmth
and strictness scores of the previous chapter: the visit-based measures, on a
0--10 scale, for the \NTierOne{} visited schools, entered per standard
deviation. The visit's teaching-practice score enters as a benchmark. The
espoused scores enter one robustness specification and the inspection-report
strictness score the national extension.

\textbf{Controls.} All specifications condition on prior attainment (mean Key
Stage~2 score of the cohort), the shares of pupils eligible for free school
meals, with English as an additional language and with special educational
needs, log school size, academy status, urban location, selective admissions,
years since the last inspection, and the school's overall inspection grade as
at August 2019, entered as dummies. The 2019 grade is used rather than the
current one because a contemporaneous grade is partly an outcome of the
culture being measured. Six visited schools converted to academies after 2019
and carry no grade of their own; for five, the predecessor school's 2019 grade
is used, and the sixth is dropped where the control is required.

\textbf{Sample.} Two exclusions apply. Schools whose statutory entry age is 13
or above (university technical colleges, studio schools and similar; 98
schools nationally, two among those visited) are excluded from every outcome
regression: Progress~8 measures progress from age 11, and these schools enrol
their pupils at 13 or 14. One visited school lacks the years-since-inspection
control. The primary estimation sample is therefore $n = 99$.
\Cref{tab:outcome_stats} and \cref{tab:control_stats} report descriptive
statistics for outcomes and controls; the correlations between the culture
scores and the controls are all below $|0.3|$ (\cref{tab:score_controls_corr}).

The visited schools lean toward stronger schools than the national
population, as the previous chapter records, so the estimates apply most
directly to a population somewhat above the national average, and the
restricted range of both culture and outcomes works against finding
associations.

% ============================================================
\section{Empirical Specification}
\label{sec:p2_spec}

The primary specification is
%
\begin{equation}
  \label{eq:main}
  Y_s = \alpha + \beta_W W_s + \beta_S S_s + \mathbf{X}_s'\gamma
        + \sum_g \delta_g \,\mathbb{1}[G_s = g] + \varepsilon_s ,
\end{equation}
%
where $Y_s$ is the two-year average Progress~8 (overall, English or Maths) of
school $s$, $W_s$ and $S_s$ are its enacted warmth and strictness scores in
standard-deviation units, $\mathbf{X}_s$ the controls listed above and $G_s$
its 2019 inspection grade. Standard errors are heteroskedasticity-consistent
(HC3). Three stages are estimated: equation~(\ref{eq:main}); the same
equation with the teaching-practice score $T_s$ in place of $W_s$ and $S_s$;
and the equation with all three. Each dimension is also entered alone. The
national extension replaces $W_s$ and $S_s$ with the inspection-report
strictness score, standardised, on all schools with an outcome, and is
estimated with and without the grade dummies.

% ============================================================
\section{Results}
\label{sec:p2_results}

\textbf{Culture predicts progress.} \Cref{tab:spec_ladder} reports the
primary estimates. In the primary specification observed warmth carries a
coefficient of $0.143$ ($p = 0.002$) and observed strictness $0.131$
($p = 0.007$) on overall Progress~8, with $R^2 = 0.62$. The table shows the
same equation under four treatments of the grade control, from requiring an
original 2019 grade ($n = 95$) to dropping the control ($n = 100$); the
coefficients move by less than $0.02$ across the four. The two dimensions
contribute with near-equal weight, as the authoritative-school framework
predicts.

\begin{table}[htbp]\centering
\def\sym#1{\ifmmode^{#1}\else\(^{#1}\)\fi}
\small
\caption{The headline gold-standard model under four treatments of the pre-COVID
inspection-grade control. Each row is the same regression of average Progress 8 on
enacted warmth ($W$) and enacted strictness ($S$) with the full control set; the rows
differ only in how schools with a missing 2019 grade are handled. Schools with a
statutory admission age of 13 or above are excluded throughout.}
\label{tab:spec_ladder}
\begin{tabular}{lccc}
\toprule
Grade-control specification & $N$ & Warmth ($W$) & Strictness ($S$) \\
\midrule
Grade as recorded (missing dropped)      &  95 & 0.155\sym{***}  & 0.128\sym{***}  \\
                                         &     & (0.045)        & (0.047)        \\
\addlinespace
Predecessor-filled grade (primary)       &  99 & 0.143\sym{***}  & 0.131\sym{***}  \\
                                         &     & (0.045)        & (0.047)        \\
\addlinespace
Missing-grade category retained          & 100 & 0.143\sym{***}  & 0.131\sym{***}  \\
                                         &     & (0.045)        & (0.047)        \\
\addlinespace
No grade control                         & 100 & 0.140\sym{***}  & 0.127\sym{***}  \\
                                         &     & (0.046)        & (0.046)        \\
\bottomrule
\end{tabular}
\begin{minipage}{\linewidth}
\smallskip
\footnotesize\textit{Notes:} Standard errors in parentheses.
\sym{*} \(p<0.10\), \sym{**} \(p<0.05\), \sym{***} \(p<0.01\).
Outcome is average Progress 8; predictors are standardised. Seven visited academies
lacked a 2019 grade under their current URN; the primary specification fills five of
them from the predecessor school's grade. Including the late-entry schools leaves the
primary estimates essentially unchanged ($W = 0.139$, $p = 0.003$;
$S = 0.133$, $p = 0.006$; $N = 100$).
\end{minipage}
\end{table}


Entered alone, each dimension carries a coefficient of about $0.22$
(\cref{tab:univariate_ws}): $0.218$ for warmth and $0.228$ for strictness.
Jointly each retains roughly 60 per cent of its solo coefficient. The two
dimensions therefore share variance, as their correlation of $0.59$ implies,
but each carries substantial signal the other does not, which is the
empirical basis for treating them as two constructs.

\begin{table}[htbp]\centering
\def\sym#1{\ifmmode^{#1}\else\(^{#1}\)\fi}
\small
\caption{Warmth and strictness entered alone and together. Each column is a Progress 8
outcome; the first two rows enter one culture dimension at a time and the final two
rows enter both jointly, always with the full control set and the predecessor-filled
grade control. The drop from the univariate to the joint coefficients shows the two
dimensions share some variance, but each keeps an independent association when the
other is held fixed.}
\label{tab:univariate_ws}
\begin{tabular}{llccc}
\toprule
Specification & Term & Overall & English & Maths \\
\midrule
$S$ only   & Strictness ($S$) & 0.195\sym{***} & 0.170\sym{***} & 0.156\sym{***} \\
           &                  & (0.033)        & (0.039)        & (0.033)        \\
\addlinespace
$W$ only   & Warmth ($W$)     & 0.206\sym{***} & 0.207\sym{***} & 0.164\sym{***} \\
           &                  & (0.035)        & (0.037)        & (0.039)        \\
\midrule
Joint      & Warmth ($W$)     & 0.132\sym{***} & 0.158\sym{***} & 0.104\sym{**}  \\
           &                  & (0.042)        & (0.041)        & (0.050)        \\
\addlinespace
Joint      & Strictness ($S$) & 0.117\sym{***} & 0.077\sym{*}   & 0.095\sym{**}  \\
           &                  & (0.039)        & (0.042)        & (0.042)        \\
\midrule
$N$        &                  & 100            & 100            & 100            \\
\bottomrule
\end{tabular}
\begin{minipage}{\linewidth}
\smallskip
\footnotesize\textit{Notes:} Standard errors in parentheses.
\sym{*} \(p<0.10\), \sym{**} \(p<0.05\), \sym{***} \(p<0.01\).
Gold-standard visited tier, late-entry schools excluded. EBacc and Open components
are reported in the chapter appendix.
\end{minipage}
\end{table}


At the component level English is warmth-led (warmth $0.167$, $p < 0.001$;
strictness $0.087$, $p = 0.083$) and Maths balanced (warmth $0.112$,
$p = 0.041$; strictness $0.110$, $p = 0.036$). The five-outcome tables in
Appendix~3A show both dimensions significant for the EBaC component and
warmth for the Open component; the appendix also shows that the strictness
score predicts entry to the EBacc pathway nationally with the grade held
constant, but that controlling for entry attenuates the strictness--EBaC
association by only 10 to 14 per cent. Most of the association is progress
within subjects rather than curriculum steering.

\textbf{Teaching quality does not account for it.} \Cref{tab:stages23_trio}
reports the second and third stages. The visit's teaching-practice score alone
predicts overall Progress~8 at $0.245$ ($p < 0.001$) but with a lower $R^2$
($0.56$) than the two culture scores. Added to the primary specification it
enters at $0.024$ ($p = 0.76$), and the culture coefficients are unchanged
($0.131$ and $0.126$). Warmth for Maths slips past the five per cent threshold
once teaching enters ($p = 0.115$), which is unsurprising given how much
variance the three visit scores share. Culture and observed instruction are
largely distinct predictors of progress; conditioning on teaching quality does
not account for the culture association, though with cross-sectional data the
alternative reading, that teaching is the noisiest of the three visit scores
and its signal is absorbed by the other two, cannot be excluded.

\begin{table}[htbp]
\centering
\small
\caption{Stages 2 and 3 under the primary specification: teaching quality as
the sole culture predictor (Stage~2), and teaching quality added to the
warmth-and-strictness model (Stage~3). Coefficients are per standard
deviation of the predictor; HC3 robust standard errors; $n = 100$. The full
five-outcome versions, estimated on the previous specification, are reported
in the appendix.}
\label{tab:stages23_trio}
\begin{tabular}{lccc}
\toprule
 & \textit{Overall P8} & \textit{English} & \textit{Maths} \\
\midrule
\multicolumn{4}{l}{\textit{Stage 2: teaching quality only}} \\
\addlinespace[2pt]
Teaching $T$ & 0.202\sym{***} & 0.184\sym{***} & 0.159\sym{***} \\
             & (0.000) & (0.000) & (0.001) \\
$R^2$        & 0.562 & 0.560 & 0.484 \\
\addlinespace[6pt]
\multicolumn{4}{l}{\textit{Stage 3: warmth, strictness and teaching}} \\
\addlinespace[2pt]
Warmth $W$   & 0.121\sym{*}  & 0.165\sym{**} & 0.098 \\
             & (0.033) & (0.006) & (0.108) \\
Strictness $S$ & 0.113\sym{**} & 0.080 & 0.093\sym{*} \\
             & (0.005) & (0.064) & (0.047) \\
Teaching $T$ & 0.020 & $-$0.010 & 0.011 \\
             & (0.758) & (0.879) & (0.872) \\
\bottomrule
\end{tabular}
\begin{minipage}{0.8\linewidth}
\vspace{4pt}\scriptsize
\textit{Notes}: $p$-values in parentheses. \sym{*} $p<0.05$, \sym{**}
$p<0.01$, \sym{***} $p<0.001$. Primary specification throughout: full
control set, predecessor-filled pre-COVID grade, late-entry schools
excluded.
\end{minipage}
\end{table}


\textbf{The dimensions add.} \Cref{fig:quadrant_ws} plots the visited schools
by their two culture scores, shaded by progress. Schools in the warm and
strict quadrant are predominantly above-average performers. Formally, the
warm-and-strict quadrant sits $0.36$ of a grade above the rest of the sample
after controls ($p < 0.001$), while the interaction between the two dimensions
is null ($0.039$, $p = 0.20$; \cref{tab:typology}). Warmth and strictness add
rather than amplify one another; the best schools have both because both help,
not because of a combination effect. Decomposing the visit scores into their
five observational blocks (\cref{tab:subscores}), the association runs
through warmth outside lessons and orderly conduct within them: out-of-lesson
warmth ($0.107$, $p = 0.017$) and in-lesson strictness ($0.118$, $p = 0.006$)
survive joint entry, and their counterparts do not.

\begin{figure}[htbp]
  \centering
  \includegraphics[width=0.85\textwidth]{fig_quadrant_ws.pdf}
  \caption{Enacted warmth and strictness for the visited schools, shaded by
  Progress~8 (two-year average). Dashed lines mark sample medians. The
  adjusted quadrant comparison is in \cref{tab:typology}.}
  \label{fig:quadrant_ws}
\end{figure}

\textbf{Espoused culture does not predict progress.} Substituting the
headteacher's espoused scores for the observed ones on the same schools
reduces the warmth coefficient to $0.036$ ($p = 0.47$); espoused strictness
retains a smaller association ($0.113$, $p = 0.019$). The measurement
asymmetry of the previous chapter reappears as a predictive one: what a
headteacher says about warmth carries no information about progress, and what
they say about behaviour systems carries some. Two related results are in
Appendix~3A. The gap between a school's espoused and enacted scores predicts
nothing beyond the enacted level; and the association is not confined to the
visited schools' own report on themselves, since parents' survey responses on
their children's happiness and safety are positively associated with progress
($0.14$ per standard deviation) with the caveat that satisfaction and results
are jointly determined.

\textbf{National extension.} The inspection-report strictness score, which
the previous chapter showed to track observed strictness at $0.40$, is
available for 3,147 schools with an outcome. On the full control set without
the grade it carries $0.166$ ($p < 0.001$, $R^2 = 0.52$;
\cref{tab:national_strictness}); with the 2019 grade added ($n = 2{,}823$) it
carries $0.167$, and no component estimate moves by more than $0.005$. A score
read from an inspection report might be thought to proxy for the inspectorate's
overall approval; that the coefficient is untouched by conditioning on the
grade itself is difficult to reconcile with that reading. Two further
text-derived results are consistent with the previous chapter's measurement
findings and are reported in \cref{tab:llm_p8_matrix}: the behaviour-policy
warmth score is significantly \emph{negatively} associated with progress, and
the website warmth and strictness scores are not associated with it at all.

\textbf{Robustness.} \Cref{tab:robustness_overall} reports the primary
estimate under alternative specifications. Dropping the grade control
($0.130$, $0.125$; $n = 102$), using the single 2023/24 outcome year ($0.165$,
$0.125$), adding the 2015/16 baseline share of pupils with social, emotional
and mental-health needs ($0.163$, $0.105$; $n = 91$), and restricting to
schools whose headteacher is confirmed unchanged since the last inspection all
leave both coefficients significant at the five per cent level or better. On
Attainment~8 rather than Progress~8 the coefficients are larger ($0.337$,
$0.218$), as expected for a measure that does not condition on intake. The
sorting concern, that strict schools have fewer pupils with social, emotional
and mental-health needs, accounts for about a tenth of the strictness
association (Appendix~3A). Within-school Progress~8 is stable across years
(mean within-school standard deviation $0.20$), which supports using culture
measured after the outcome window as a proxy for culture during it; the
appendix also shows that the 2024/25 pseudo-progress measure, the one cohort
measured contemporaneously with the fieldwork, gives coefficients within
$0.04$ of the primary ones.

\textbf{Implications.} Two implications bear on current policy debate and
follow from the pattern of results rather than any single coefficient. The
behaviour debate treats discipline as the active ingredient of school
improvement and relational culture as secondary; these estimates give the two
equal weight, show no trade-off between them, and place the warm-and-strict
schools a third of a grade ahead. And the association runs through relational
quality outside lessons and orderly conduct within them, which locates the
relational work that matters outside the classroom, where observation and
training effort are least concentrated. Both implications are conditional on
an association being at least partly causal, which the design supports but
does not establish.

% ============================================================
\section{Conclusion}
\label{sec:p2_conclusion}

Observed school culture predicts pupil progress. Across \NTierOne{} visited
English secondary schools, a standard deviation of observed warmth is
associated with $0.14$ of a grade higher Progress~8 and a standard deviation
of observed strictness with $0.13$, both significant at the one per cent
level and both stable across every reasonable specification. The two
dimensions contribute equally and additively; the schools that combine them
outperform the rest by a third of a grade. Observed teaching quality does not
account for the association. The headteacher's own account of the school
predicts nothing on warmth, and the strictness score read from inspection
reports reproduces the strictness result nationally with the inspection grade
held constant.

The estimates are associations, strengthened by a value-added outcome,
temporal ordering, a rich control set and stability across specifications, but
not identified: a persistent characteristic that produces both a school's
culture and its results cannot be excluded with cross-sectional data. The
visited sample is small and leans toward stronger schools; culture was
measured after the outcome window for most schools; and one visit is a noisy
reading of a school. Within those limits, the results locate the relational
and disciplinary features of a school that travel with its pupils' progress,
and they show that one of the two, strictness, is legible at national scale
in the public record.
